\documentclass{ximera}
\title{Describing Simple Harmonic Motion}

\newcommand{\pskip}{\vskip 0.1 in}

\begin{document}
\begin{abstract}
Using the sine and cosine functions to describe simple harmonic motion.
\end{abstract}
\maketitle





\begin{example}  \label{E8bdsdsf}
Suppose that over the course of a 24-hour period, from midnight October 29 to midnight October 30, the depth of the water at the Edmonds Pier varies in simple harmonic motion about its mean. Suppose also that a high tide of 21 feet occurs at 2:00am and the following low tide of 5 feet occurs at 8:00am. 

We'll ask two questions:

\begin{enumerate}
\item Use circular interpolation with the protractor below to approximate the depth of the water at 3:00am on October 29.

\item Use circular interpolation with the protractor below to approximate the depth of the water at 5:30pm on October 29.

%\item %Use the circular protractor to approximate the first two times after midnight when the water is $18$ feet deep.

Explain your reasoning and include screenshots to help with your explanations. 
\begin{freeResponse}
\end{freeResponse}

\begin{onlineOnly}
    \begin{center}
\desmos{0wxwmkzvky}{900}{600}
\end{center}
\end{onlineOnly}

\href{https://www.desmos.com/calculator/0wxwmkzvky}{142: Circular Interpolation Tides}

%\emph{Keep in mind that high tide occurs at 2:00am and low tide at 8:00am.}

\end{enumerate}
\end{example}




\section{Circular Interpolation using the Sine Function}



\begin{example} \label{Efoe34r}
This is a contintuation of Example 1. Our goal now is to find a function
\[
     h = f(t) \, , \, 0\leq t \leq 24,
\]
that expresses the depth of the water (in feet) at the Edmonds pier in terms of the number of hours past midnight, October 29.

Here are the steps.

\begin{enumerate}

\item  First we'll use the information above to sketch by hand a graph of the function $f$. Label the axes with the appropriate variable names and units. Label the coordinates of four key points on the graph.

Then activate the folder \emph{Graph} in Line 2 of the worksheet below to check your graph.

\begin{onlineOnly}
    \begin{center}
\desmos{jqmcce5as4}{900}{600}
\end{center}
\end{onlineOnly}

\href{https://www.desmos.com/calculator/jqmcce5as4}{142: Edmonds Pier 2a}.

\item Next we'll compute the mean (average) depth of the water over the course of a $12$-hour and the maximum deviation from this mean. Including units in our computation, the mean depth is 
\[
    h_{avg} = 0.5 ( 21 \text{feet } + 5 \text{feet}) = 13 \text{ feet} . 
\]
And the maximum deviation from the mean is
\[
    r = 0.5 ( 21 \text{feet } - 5 \text{feet}) = 8 \text{ feet} ..
\]

Activate the \emph{Amplitude} folder on Line 8 of the above demonstration to draw the horizontal line showing the average depth. Note this line is labeled with its equation. There is also a vertical line that shows the the maximum deviation from the mean.

\item Next we'll draw a circle on a new $x,y$-coordinate system centered at any point with $h$-coordinate $h_{\text{avg}}=13$ feet equal to the mean depth and radius $r=8$ feet equal to the maximum deviation from the mean.

Activate the folders \emph{Uniform circular motion} and \emph{Polar angle} on Lines 21 and 29 to see the circle and coordinate system.

\item Now the key point: The constant-speed motion of point $Q$ around this circle drives the variation of the water's depth. As point $Q$ turns about the circle's center at a constant rate, the point $P$ moves along the graph of the function $h=f(t)$.

To start, we'll forget about time and express the depth of the water (in feet) in terms of the polar angle $\theta$ of $Q$. There are two steps.

First we express the $y$-coordinate of $Q$ in terms of $\theta$. This is quick. Since the circle has radius $8$ feet, $Q$ has $y$-coordinate (in feet),
\[
     y = 8\sin\theta .
\]
Then since the circle is centered at the mean depth $h_{\text{avg}}=13$ feet, the depth as a function of the polar angle is
\[
     h =g(\theta) = 13 +  8\sin\theta \, , \, 0\leq \theta \leq 4\pi. 
\]

\item Finally, to express the depth as a function of time, we need only find a function
\[
     \theta = a(t) \, , \, 0\leq t \leq 24,
\]
that expresses the polar angle (in radians) in terms of the number of hours past midnight.

We do this in the usual way.  First we compute the rotation rate of the driving point $Q$ about circle's center. The polar angle at the high tide at 2:00am is
\[
   a(2) = \pi/2 \text{ radians}.
\]
At the following low tide (8am), the polar angle is
\[
  a(8) = 3\pi/2 \text{ radians}.
\]
So $Q$ rotates about the origin of the $xy$-coordinate system at the constant rate of

The rotation rate is
\begin{align*}
    \frac{\Delta \theta}{\Delta t} &= \frac{\left( \frac{3\pi}{2} - \frac{\pi}{2}\right)\text{ radians}}{(8-2)\text{ hrs}} \\
                                              &= \frac{\pi}{6} \text{ radians/hr}.
\end{align*}

Since $a(2) = \pi/2$, the polar angle function is
\[
   \theta = a(t) = \frac{\pi}{2} + \frac{\pi}{6} \left(  t-2 \right) .
\]
Substituting this expression for the polar angle in the function
\[
 h = g(\theta) = 13 + 8\sin\theta \, , \, 0\leq \theta \leq 24,
\]
we get
\begin{align*}
 h  &= f(t) \\
     &= g(a(t)) \\
     &= g\left( \frac{\pi}{2} + \frac{\pi}{6} \left(  t-2 \right) \right) \\
     &= 13 + 8 \sin\left( \frac{\pi}{2} + \frac{\pi}{6} \left(  t-2 \right) \right) \, , \, 0\leq \theta \leq 24,
\end{align*}
for the function that expresses the water's depth (in feet) in terms of the number of hours past midnight, October 29.

\item To check that our function is correct, we'll use the given information that at 2am ($t=2$) the depth is $21$ feet and  $5$ feet at 8:00am.

Substituting $t=2$, we get
\begin{align*}
   f(2)   & = 13 + 8 \sin \left(\frac{\pi}{2}  \frac{\pi}{6} \left( 2-2 \right) \right) \\ 
           & = 13 + 8 \sin (\frac{\pi}{2})  \\ 
           & = 13 + 8 \\
           & = 21 .
\end{align*}

Substituting $t=8$, we get
\begin{align*}
   f(8)   & = 13 + 8 \sin \left(\frac{\pi}{2}  \frac{\pi}{6} \left( 8-2 \right) \right) \\ 
           & = 13 + 8 \sin \left( \frac{3\pi}{2} \right)  \\ 
           & = 13 + 8(-1) \\
           & = 5 .
\end{align*}
These check out.

\end{enumerate}
\end{example}


\section{Circular Interpolation Using the Coine Function}
\begin{example} \label{Ex:lfffme3f}
This is a continuation of Example 1, but now we use the cosine function instead to find an expression for the function
\[
     h = f(t) \, , \, 0\leq t \leq 24,
\]
that expresses the depth of the water (in feet) at the Edmonds pier in terms of the number of hours past midnight, October 29. This way is a bit easier.

The first change is to rotate the axes of the $x,y$-coordinate system so that the $x$-axis points vertically upward, in the same direction as the $h$-axis. This, while not strictly necessary, makes it easier to see the relationship between the $x$-coordinate of the driving point and the depth of the water. But as always, we measure the polar angle counterclockwise from the positive $x$-axis.

\begin{onlineOnly}
    \begin{center}
\desmos{caylccjijk}{900}{600}
\end{center}
\end{onlineOnly}

\href{https://www.desmos.com/calculator/caylccjijk}{142: Edmonds Pier 2c}.

To find the depth of the water as a function of the polar angle, we start with the mean depth of 13 feet and add to that the $8$-foot oscillation, this time using the cosine function. So (omitting the domain) the depth as a function of the polar angle is
\[
    h = g(\theta) = 13 + 8\cos\theta .
\]
The rotation rate of the driving point is still $\pi/6$ radians/hour. All that's changed is the values of the polar angle. At high tide ($t=2)$ the polar angle is now $a(2) = 0$. So the polar angle function is
\[
    \theta = a(t) = \frac{\pi}{6} \left( t - 2 \right) \, , \, 0\leq t \leq 24.
\]  
And the depth (in feet) as a function of the number of hours past midnight is
\begin{align*}
     h &= f(t)  \\
         & = g(a(t)) \\
         &= g\left(\frac{\pi}{6} \left( t - 2 \right) \right)  \\
         &= 13 + 8\cos \left(  \frac{\pi}{6} \left( t - 2 \right) \right) \, , \, 0\leq t \leq 24.
\end{align*}

\end{example}








\end{document}